(i) Subtracting consecutive equations of \eqref{eq:run} eliminates $a_{U}$ and gives the linear system $b_{12}x + \Dl_{12}y = D_{12}$, $b_{23}x + \Dl_{23}y = D_{23}$ in $(x, y)$, whose determinant is $b_{12}\Dl_{23} - b_{23}\Dl_{12}$; Cramer's rule gives the displayed $x$ and $y$, and the third equation gives $a_{U}$. Dividing, $f = (b_{12}D_{23} - b_{23}D_{12})/(D_{12}\Dl_{23} - D_{23}\Dl_{12})$, and dividing numerator and denominator by $D_{23}$ gives the form in $R$. When the determinant vanishes the two equations are proportional on the left, so a solution exists only if $D_{12}/D_{23} = b_{12}/b_{23}$. The script substitutes the closed forms back into \eqref{eq:run} for every spectrum of Proposition \claimref{standard-model.table-exact} and checks the identity for $f$.

(ii) $\Dl_{12}$ and $\Dl_{23}$ are unchanged by the shift, so the $(x, y)$ system and its solution are unchanged; $a_{U} = a_{3} - b_{3}x - (\Dl_{3} + c)y$ drops by $cy$. The script checks this for three values of $c$ on every spectrum. Lemma \claimref{standard-model.complete-multiplets-shift-all-three-coefficients} supplies the content with $\Dl \propto (1,1,1)$.

(iii) For $t \in [0, 1]$ set $r(t) = (b_{12} + t\Dl_{12})/(b_{23} + t\Dl_{23})$, so $r(0) = r_{0}$ and $r(1) = r_{\Dl}$. The denominator is positive at both ends and linear, hence positive on $[0, 1]$, so $r$ is a M\"obius function of $t$ without a pole on the interval, with derivative $(\Dl_{12}b_{23} - \Dl_{23}b_{12})/(b_{23} + t\Dl_{23})^{2}$ of constant sign: $r$ is strictly monotone on $[0, 1]$, or constant when the determinant vanishes. Now the two equations of (i) say precisely $D_{12} = x(b_{12} + f\Dl_{12})$ and $D_{23} = x(b_{23} + f\Dl_{23})$, so $R = r(f)$ and $x = D_{23}/(b_{23} + f\Dl_{23})$. Hence a solution with $f \in [0, 1]$ exists iff $R$ is a value of $r$ on $[0, 1]$, iff $R$ lies between $r(0)$ and $r(1)$ by monotonicity and continuity; and then $x > 0$ since $D_{23} > 0$ and the denominator is positive. The script checks the equivalence and the monotonicity (sampled at ninths, in exact arithmetic) on every spectrum of the table.
